%\vspace{-5pt}
\begin{abstract}
%\vspace{-10pt}
We present Habitat, a platform for research in embodied artificial intelligence (AI).
Habitat enables training embodied agents (virtual robots) in highly efficient photorealistic 3D simulation.
Specifically, Habitat consists of: \\
(i) \habitatsim: a flexible, high-performance 3D simulator with configurable agents, sensors, and generic 3D dataset handling. \habitatsim is fast -- when rendering a scene from Matterport3D, it achieves several thousand frames per second (fps) running single-threaded, and can reach over $10@000$ fps multi-process on a single GPU. \\ 
(ii) \habitatapi: a modular high-level library for end-to-end development of embodied AI algorithms -- defining tasks (\eg navigation, instruction following, question answering), configuring, training, and benchmarking embodied agents.
%\end{asparaenum}

These large-scale engineering contributions enable us to answer scientific questions requiring experiments that were till now impracticable or `merely' impractical.
Specifically, in the context of point-goal navigation:
(1) we revisit the comparison between learning and SLAM approaches from two recent works~\cite{mishkin2019benchmarking, kojima2019learn} and find evidence for the \textbf{\emph{opposite conclusion}} --
that learning outperforms SLAM if scaled to an order of magnitude 
more experience than previous investigations, and
(2) we conduct the first cross-dataset generalization experiments 
$\{$train, test$\} \times \{$Matterport3D, Gibson$\}$ for multiple sensors $\{$blind, RGB, RGBD, D$\}$ and find that only agents with depth (D) sensors generalize across datasets.
We hope that our open-source platform and these findings will advance research in embodied AI.

\end{abstract}